\documentclass[11pt,letterpaper]{article}
\PassOptionsToPackage{hyphens}{url}
\usepackage[margin=.83in,top=.8in,bottom=.8in,headheight=15pt]{geometry}
\usepackage[T1]{fontenc}
\usepackage{lmodern,amsmath,amssymb,mathtools,bm,booktabs,array,graphicx,tikz}
\usetikzlibrary{arrows.meta,positioning,calc,decorations.pathreplacing,patterns}
\definecolor{ink}{HTML}{17344D}\definecolor{accent}{HTML}{087E8B}\definecolor{pale}{HTML}{EDF5F5}\definecolor{gold}{HTML}{AD6B19}
\providecommand{\dd}{\,\mathrm{d}}\providecommand{\R}{\mathbb{R}}
\providecommand{\E}{\mathbb{E}}\providecommand{\Var}{\operatorname{Var}}
\providecommand{\Cov}{\operatorname{Cov}}\providecommand{\diag}{\operatorname{diag}}
\providecommand{\norm}[1]{\left\lVert#1\right\rVert}
\providecommand{\abs}[1]{\left|#1\right|}
\usepackage{microtype,hyperref,fancyhdr,enumitem,amsthm,sectsty,needspace}
\allsectionsfont{\raggedright}
\hypersetup{colorlinks=true,linkcolor=ink,urlcolor=accent,pdfauthor={Zan Ahmad},pdftitle={A mixed chamber: differential equations and numerical stability / Review}}
\setlength{\parindent}{0pt}\setlength{\parskip}{6pt}\linespread{1.07}
\setlist{itemsep=3pt,topsep=4pt,leftmargin=*}
\emergencystretch=2em\clubpenalty=10000\widowpenalty=10000
\pagestyle{fancy}\fancyhf{}\fancyhead[L]{\small Mathematical Modeling}\fancyhead[R]{\small Lecture 3}
\fancyfoot[L]{\footnotesize Zan Ahmad}\fancyfoot[R]{\footnotesize\thepage}\renewcommand{\headrulewidth}{.3pt}
\newtheorem{proposition}{Proposition}\newtheorem{theorem}{Theorem}\newtheorem{lemma}{Lemma}
\theoremstyle{definition}\newtheorem{definition}{Definition}\newtheorem{example}{Example}
\numberwithin{equation}{section}
\begin{document}
\section*{A mixed chamber: differential equations and numerical stability --- review sheet}
Use this after attempting the worked examples and practice. Explain each claim in words before calculating.

\section*{Key ideas}
\begin{itemize}
\item Amount and rate have different units; change =\allowbreak{} rate $\times$ elapsed time.
\item Equilibrium balances input and removal; r/k only applies when k>0.
\item Explicit Euler uses the old state in the rate.
\item Compare solutions at the same physical time.
\item Conservation, positivity, stability and accuracy are different properties.
\end{itemize}
\section*{Common errors and repairs}
\textbf{Adding mg/min directly to mg.} Multiply the entire net rate by a time interval.

\textbf{Calling a conserved but negative amount acceptable.} Check physical admissibility separately from bookkeeping.

\textbf{Calling the circulation code explicit Euler.} This chamber is an explicit teaching example; the source circulation solver is implicit with a valve-consistency loop.

\section*{Retrieve without your notes}
\begin{itemize}
\item What does 1/k measure?
\item Why might Euler fill too quickly from an empty chamber?
\item What happens if k=\allowbreak{}0?
\end{itemize}
\section*{Readiness check}
\begin{itemize}
\item Distinguish an amount from a rate and convert a rate into a change over a time interval.
\item Build an explicit Euler update from a conservation statement.
\item Predict equilibrium and compare a numerical trajectory with an exact reference.
\item Recognize that conservation alone does not guarantee positivity or accuracy.
\end{itemize}
\end{document}
